\chapter{The FRI Protocol}
\label{ch:fri}
 
This chapter presents the Fast Reed-Solomon Interactive Oracle Proof of
Proximity, the protocol at the core of the system implemented in
this thesis. We begin by stating the problem FRI solves, then describe
the algebraic structure of the evaluation domain, and finally the
folding operation that drives the protocol.

\section{Overview and Problem Statement}
\label{sec:fri_overview}

 
Recall from \cref{def:rs_code} that, for a subset $S \subseteq \F$ and
a rate parameter $\rho \in (0,1)$, the Reed-Solomon code
$\RS[\F, S, \rho]$ consists of the evaluations on $S$ of polynomials of
degree less than $\rho|S|$. The problem FRI addresses is the following.
 
\begin{definition}[Rate Proximity Testing (RPT) {\cite[Section~2.1]{BBHR18}}]
\label{def:rpt}
A verifier is given oracle access to a function $f : S \to \F$, together
with auxiliary information supplied by a prover. The verifier must
distinguish, with high probability and using only a small number of
queries, between the case that $f \in \RS[\F, S, \rho]$ and the case
that $f$ is far from every member of $\RS[\F, S, \rho]$ in relative
Hamming distance (\cref{def:hamming}).
\end{definition}
 
RPT is a special case of the low-degree testing problem, and
Ben-Sasson et al.\ identify it as a major bottleneck for transparent
proof systems~\cite[Section~2.1]{BBHR18}.
 
A naive solution is instructive. The verifier could ask the prover to
open $f$ at $\rho|S|$ positions, interpolate the unique polynomial of
degree less than $\rho|S|$ through them, and check that it agrees with
$f$ elsewhere. This fails on two counts: the number of openings is
proportional to the degree bound, which can reach $2^{20}$ or more in
practical instances, and every opening carries a Merkle authentication
path of $\log|S|$ hash values (\cref{sec:merkle_trees}), so the proof
size becomes prohibitive.
 
FRI takes a different route. Instead of testing the degree directly,
it \emph{reduces} it: the prover repeatedly halves the degree of the
committed polynomial using randomness supplied by the verifier, until
the degree is small enough to send the remaining polynomial directly.
Each reduction step is committed, and the verifier spot-checks the
consistency of the steps at a small number of random positions. If the
original function is far from the code, then with high probability at
least one reduction is inconsistent, and the verifier detects it.
 
This yields a protocol with prover arithmetic complexity strictly
linear in $|S|$ and verifier arithmetic complexity strictly
logarithmic~\cite[Section~2.1]{BBHR18}. In the terminology of
\cref{sec:iop}, FRI is an interactive oracle proof of proximity (IOPP)
for Reed-Solomon codes: the prover's messages are committed
codewords which the verifier queries at a bounded number of positions,
rather than reading in full.
 
\begin{remark}
\label{rem:proximity_not_membership}
FRI establishes \emph{proximity}, not membership. A successful
verification does not certify that the committed vector is exactly a
Reed-Solomon codeword, only that it is close to one in relative
Hamming distance. This distinction is central to the soundness
analysis of the FRI protocol.
\end{remark}  % citare soundess proof
 

 
% ==================================================================
\section{The Evaluation Domain}
\label{sec:fri_domain}
% ==================================================================
 
The folding step of FRI relies on specific algebraic structure in the
evaluation domain. We state the requirements first, then give the
construction used in our implementation.
 
 
\subsection{Requirements}
\label{subsec:domain_requirements}
 
The reduction from a polynomial $f$ to a polynomial of half the degree
is carried out by pairing the evaluations of $f$ at $x$ and at $-x$. %(\cref{sec:fri_folding})
For this to be possible at every position,
the domain must be closed under negation. Moreover, the folded
polynomial is committed not on $D$ but on the squared domain
$D^{(2)} := \{x^2 : x \in D\}$, and since the folding step is applied
iteratively, $D^{(2)}$ must be closed under negation as well, and so
on for every subsequent layer~\cite[Section~3.9.1]{ethSTARK}.
 
We therefore require a chain of domains
%
\[
D = D_0, \quad D_{i+1} = D_i^{(2)}, \quad i = 0, 1, \ldots
\]
%
such that every $D_i$ is closed under negation and
$|D_{i+1}| = |D_i|/2$. These requirements are met by taking $D$ to be
a coset of a multiplicative subgroup of order
$2^k$~\cite[Section~3.9.1]{ethSTARK}.
 
A coset is used rather than the subgroup itself so that the evaluation
domain is disjoint from the trace domain on which the polynomial being
committed is defined~\cite[Section~3.4]{ethSTARK}. Disjointness is
required in the full STARK protocol, where constraint polynomials are
expressed as rational functions whose denominators vanish on the trace
domain; evaluating them on a disjoint domain keeps those denominators
nonzero.
 
 
\subsection{Construction}
\label{subsec:domain_construction}

We build the domain from the roots of unity of \cref{subsec:roots_of_unity}.
Let $n$ be the number of coefficients of the polynomial to be committed,
$b$ the blowup factor, and $N = b \cdot n = 2^k$ the size of the
evaluation domain. Two elements are fixed:
%
\begin{equation}
\label{eq:domain_generators}
\omega = \gamma^{(p-1)/2^{k}}, \qquad
g = \gamma^{(p-1)/2^{k+1}},
\end{equation}
%
where $\gamma$ is a fixed multiplicative generator of $\Fp^*$. Here
$\omega$ is a primitive $2^k$-th root of unity: its powers
$\langle\omega\rangle = \{1, \omega, \ldots, \omega^{N-1}\}$ form a
cyclic group of $N$ equally spaced points. The element $g$ is a
primitive $2^{k+1}$-th root of unity, one ``half-step'' finer, and it
plays the role of an \emph{offset} that rotates the whole group. The
two are linked by the invariant
%
\begin{equation}
\label{eq:g_squared}
g^2 = \omega,
\end{equation}
%
which, as we show below, is exactly what makes the domain reproduce
itself under squaring.

The evaluation domain is the coset
%
\begin{equation}
\label{eq:domain}
D = g \cdot \langle \omega \rangle
  = \{\, g\omega^j : j = 0, \ldots, N-1 \,\},
\end{equation}
%
i.e.\ the group $\langle\omega\rangle$ shifted by the offset $g$. The
codeword committed by the prover is $\bigl(f(g\omega^j)\bigr)_{j=0}^{N-1}$,
computed by scaling the $i$-th coefficient of $f$ by $g^i$ and applying
the NTT (\cref{subsec:ntt,subsec:lde}). In our implementation $\omega$
and $g$ are two-power roots of unity obtained from the fixed generator
$\gamma$, and the codeword is obtained by the scaled NTT of
\cref{subsec:lde}.

Three properties follow, and together they justify the folding step.
 
\begin{proposition}
\label{prop:domain}
With $D$ as in \cref{eq:domain} and $N = 2^k$, $k \geq 1$:
\begin{enumerate}
  \item $\omega^{N/2} = -1$;
  \item $D$ is closed under negation: $-\,g\omega^{j} = g\omega^{j+N/2}$
    for $j < N/2$;
  \item $D^{(2)} = \omega \cdot \langle\omega^2\rangle$, a coset of the
    subgroup of order $N/2$; in particular $|D^{(2)}| = N/2$.
\end{enumerate}
\end{proposition}

\begin{proof}
\emph{(1)} The element $\omega^{N/2}$ squares to
$\omega^N = 1$, so it is a square root of unity. It cannot equal $1$,
since $\omega$ has order exactly $N$; as $\Fp$ is a field, the only
square roots of unity are $\pm 1$, hence $\omega^{N/2} = -1$.

\emph{(2)} Using (1),
$g\omega^{j+N/2} = g\omega^{j}\cdot\omega^{N/2} = -\,g\omega^{j}$, and
the index $j+N/2$ lies in $\{0,\dots,N-1\}$, so the negated element is
again in $D$.

\emph{(3)} Squaring an element of $D$ and applying \cref{eq:g_squared}
gives $(g\omega^{j})^2 = g^2\omega^{2j} = \omega^{2j+1}$, so
$D^{(2)} = \{\omega^{2j+1} : j = 0,\dots,N-1\}$, the subgroup
$\langle\omega^2\rangle$ shifted by $\omega$. By (2) the map
$x \mapsto x^2$ identifies $x$ with $-x$, so it is two-to-one on $D$
and $|D^{(2)}| = N/2$.
\end{proof}

Property~(3) is what makes the construction \emph{self-reproducing}:
the squared domain is again a coset of a two-power subgroup, of the
same shape as $D$ but half the size, now described by the pair
$(\omega, \omega^2)$ in place of $(g, \omega)$. A single FRI round
therefore updates the domain simply by squaring both parameters,
%
\begin{equation}
\label{eq:domain_update}
g \leftarrow g^2, \qquad \omega \leftarrow \omega^2,
\end{equation}
%
that is, the two domain parameters are squared at the end of each
layer. The requirements of
\cref{subsec:domain_requirements} are preserved at every level, as long
as the subgroup order remains even.

\begin{example}
\label{ex:domain_f17}
Take $\F_{17}$, whose multiplicative group has order $16 = 2^4$, with
generator $\gamma = 3$. For $N = 8$ we get $\omega = 3^2 = 9$ and
$g = 3^1 = 3$, so $g^2 = 9 = \omega$ as required by
\cref{eq:g_squared}. The evaluation domain is
%
\[
D = \{3, 10, 5, 11, 14, 7, 12, 6\}.
\]
%
Its elements pair up as $(3,14)$, $(10,7)$, $(5,12)$, $(11,6)$, each
pair summing to $0 \bmod 17$, this is $-\,g\omega^{j} = g\omega^{j+4}$
in action. Squaring gives $D^{(2)} = \{9, 15, 8, 2\}$, of size $4$,
which is exactly the coset
$\omega\cdot\langle\omega^2\rangle = 9\cdot\langle 13\rangle$ built from
the updated parameters.
\end{example}


\section{The Folding Operation}
\label{sec:fri_folding}

Folding is the algebraic core of FRI. In one step it takes a polynomial
$f$ of degree less than $d$, given by its values on the domain $D$, and
produces a new polynomial of degree less than $d/2$, given by its values
on the smaller domain $D^{(2)}$. The reduction uses a single random
challenge and no interpolation: the values of the new polynomial are a
simple combination of the values of the old one. Repeating this step
reduces the degree by half each time; the prover stops once the
polynomial is small enough to send directly, and transmits its remaining
coefficients in the clear.

The section is organised around the three questions that folding
answers: how a polynomial splits into two halves
(\cref{subsec:even_odd}), how those halves are recovered from evaluations
alone (\cref{subsec:recovering}), and how the challenge combines them
into the folded polynomial (\cref{subsec:the_fold}). The two polynomials
$f_e$ and $f_o$ introduced next appear throughout the section in three
guises. They are first defined as polynomials, then evaluated at a point,
then recombined, but they are always the same two objects.


\subsection{Even-Odd Decomposition}
\label{subsec:even_odd}

The first step is to split $f$ into the part carried by its even-indexed
coefficients and the part carried by its odd-indexed ones. Grouping the
coefficients this way lets us write $f$ in terms of two polynomials in
$X^2$, each of roughly half the degree.

\begin{definition}[Even-Odd Decomposition]
\label{def:even_odd}
Let $f(X) = \sum_{i=0}^{d-1} c_i X^i \in \Fp[X]$. Define
%
\[
f_{e}(Y) = \sum_{i \text{ even}} c_i \, Y^{i/2},
\qquad
f_{o}(Y) = \sum_{i \text{ odd}} c_i \, Y^{(i-1)/2}.
\]
%
Then $f$ decomposes uniquely as
%
\begin{equation}
\label{eq:even_odd}
f(X) = f_{e}(X^2) + X \cdot f_{o}(X^2),
\end{equation}
%
and both $f_e$ and $f_o$ have degree less than $\lceil d/2 \rceil$.
\end{definition}

In words: $f_e$ collects the even coefficients $c_0, c_2, c_4, \ldots$
and $f_o$ the odd ones $c_1, c_3, c_5, \ldots$, each repackaged as a
polynomial in $Y = X^2$. Equation~\eqref{eq:even_odd} just puts them back
together, with the extra factor $X$ accounting for the odd powers. Each
half keeps at most $d/2$ of the original $d$ coefficients, which is why
its degree is halved. This is the same split that drives the recursive
step of the NTT (\cref{subsec:ntt}), and it is the decomposition FRI uses
in the commit phase.


\subsection{Recovering the Components from a Pair}
\label{subsec:recovering}

There is a catch: the prover does not hold $f$ as a list of coefficients,
but as a list of values on $D$. So the values of $f_e$ and $f_o$ must be
recovered from the values of $f$ alone. The trick is to evaluate $f$ at a
point $x$ and at its negative $-x$. Because $(-x)^2 = x^2$, both
evaluations land on the same point $x^2$ of the squared domain, and the
odd part flips sign:
%
\begin{equation}
\label{eq:fold_system}
\begin{aligned}
f(x)  &= f_{e}(x^2) + x \cdot f_{o}(x^2), \\
f(-x) &= f_{e}(x^2) - x \cdot f_{o}(x^2).
\end{aligned}
\end{equation}
%
These are two linear equations in the two unknowns $f_{e}(x^2)$ and
$f_{o}(x^2)$. Adding them isolates the even part; subtracting them
isolates the odd part. Since $x \neq 0$ for every $x \in D$, we can
divide and obtain
%
\begin{equation}
\label{eq:fold_solve}
f_{e}(x^2) = \frac{f(x) + f(-x)}{2},
\qquad
f_{o}(x^2) = \frac{f(x) - f(-x)}{2x}.
\end{equation}
%
So a single pair of values $\bigl(f(x), f(-x)\bigr)$ determines both
components at the point $x^2 \in D^{(2)}$. By \cref{prop:domain} that pair
is stored in the codeword at indices $i$ and $i + N/2$, so the prover
finds it by reading two entries that sit exactly half the list apart.


\subsection{The Fold}
\label{subsec:the_fold}

The final step mixes the two halves using the verifier's random
challenge $\alpha$. The mix collapses $f_e$ and $f_o$ into one polynomial
of half the degree.

\begin{definition}[FRI Fold]
\label{def:fold}
Let $f$ have degree less than $d$ and let $\alpha$ be a challenge sampled
from a field $\mathbb{K} \supseteq \Fp$. The \emph{fold} of $f$ with
respect to $\alpha$ is
%
\begin{equation}
\label{eq:fold}
f^{(\alpha)}(Y) := f_{e}(Y) + \alpha \cdot f_{o}(Y) ,
\end{equation}
%
a polynomial over $\mathbb{K}$ of degree less than $\lceil d/2 \rceil$.
\end{definition}

Putting \cref{eq:fold} together with \cref{eq:fold_solve}, the values of
$f^{(\alpha)}$ on $D^{(2)}$ come straight from the values of $f$ on $D$,
with no interpolation anywhere:
%
\begin{equation}
\label{eq:fold_evaluations}
f^{(\alpha)}(x^2)
  = \frac{f(x) + f(-x)}{2}
  + \alpha \cdot \frac{f(x) - f(-x)}{2x} .
\end{equation}
%
This is exactly what our implementation does: for each $i < N/2$ it reads
the two entries at indices $i$ and $i + N/2$, applies
\cref{eq:fold_evaluations} with $x = g\omega^{i}$, and writes out a
codeword of length $N/2$ on the domain $D^{(2)}$ given by the updated
parameters of \cref{eq:domain_update}.

One round thus halves two things at once: the degree bound, by
\cref{def:fold}, and the domain size, by \cref{prop:domain}. The rate
$\rho$ stays the same, so the folded codeword is again a Reed-Solomon
codeword of the same rate, which is what lets the next round repeat the
identical step. The prover keeps folding until the degree is small enough to send the
remaining polynomial directly.

% alpha is sampled from F_{p^4} by the challenger; from round 0 onwards
% the codeword entries are extension field elements.
\begin{remark}[Arithmetic in the extension field]
\label{rem:ext_field_fold}
The initial codeword has entries in $\Fp$, but the challenge $\alpha$ is
sampled from the quartic extension $\Fpk$ (\cref{subsec:field_ext}).
Because \cref{eq:fold} multiplies $f_o$ by $\alpha$, the folded
polynomial already has coefficients in $\Fpk$, and so does every layer
after it. In our implementation this switch happens at the very first
fold: the initial codeword holds base-field elements, every later layer
holds extension-field elements, and the Merkle leaves grow accordingly
from four to sixteen bytes.
\end{remark}

\begin{example}
\label{ex:fold_f17}
We continue \cref{ex:domain_f17}. Let
$f(X) = 1 + 2X + 3X^2 + 4X^3$ over $\F_{17}$, so $d = 4$, $N = 8$, and
$\rho = 1/2$. Evaluating $f$ on $D = \{3, 10, 5, 11, 14, 7, 12, 6\}$
gives the codeword
%
\[
\bigl(f(x)\bigr)_{x \in D} = (6, 3, 8, 15, 16, 4, 8, 16).
\]
%
Splitting $f$ as in \cref{def:even_odd} gives $f_e(Y) = 1 + 3Y$ and
$f_o(Y) = 2 + 4Y$. With the challenge $\alpha = 5$, the fold is
%
\[
f^{(5)}(Y) = (1 + 3Y) + 5\,(2 + 4Y) = 11 + 6Y ,
\]
%
of degree $1 < 2$, exactly as \cref{def:fold} predicts. Applying
\cref{eq:fold_evaluations} to the four pairs $\bigl(f(x), f(-x)\bigr)$
gives the codeword $(14, 16, 8, 6)$ on $D^{(2)} = \{9, 15, 8, 2\}$, which
is precisely the evaluation of $11 + 6Y$ on that domain. Both the degree
bound and the domain size have halved, and the rate is still
$2/4 = 1/2$.
\end{example}

\section{The Protocol}
\label{sec:fri_protocol}

The folding operation of \cref{sec:fri_folding} is the engine; the
protocol is what wraps it into an interaction between prover and
verifier. It has two phases. In the \emph{commit phase} the prover folds
the polynomial down to a short final layer, committing to each
intermediate codeword with a Merkle tree so that it cannot change its
mind later. In the \emph{query phase} the verifier picks a few random
positions and asks the prover to reveal the values needed to re-check
that each fold was done correctly. If the prover cheated, that is, if the
starting codeword was not close to a low-degree polynomial, then at least
one of these checks fails with high probability.

Throughout, $D_0 = D$ is the initial domain of size $N$, and
$D_{r}$ is the domain after $r$ folds, of size $N/2^r$.


\subsection{Commit Phase}
\label{subsec:commit_phase}

The prover starts from the coefficients of $f$, extends them to a
codeword on $D$ by the low-degree extension (\cref{subsec:lde}), and then
repeats the same three actions at every round: commit, receive a
challenge, fold.

\begin{enumerate}
  \item \textbf{Commit.} Send the oracle for the current codeword. This
    binds the prover to that exact codeword: every later query must be
    answered consistently with it.
  \item \textbf{Challenge.} Obtain the folding challenge $\alpha_r$ from
    the verifier (once made non-interactive by the BCS transformation of
    \cref{sec:bcs}, this challenge is derived from the transcript).
  \item \textbf{Fold.} Apply \cref{eq:fold_evaluations} with $\alpha_r$
    to halve the codeword, and square the two domain parameters
    (\cref{eq:domain_update}) to move to $D_{r+1}$.
\end{enumerate}

The process stops once the codeword has shrunk to a short final layer,
which the prover sends directly. The number of rounds is
$\log_2 N$ down to a constant, or $\log_2(N/\ell)$ if the prover stops
at a final layer of length $\ell$.

In the IOP the prover sends, at each round, the oracle for the current
codeword, and the verifier replies with the folding challenge; the
concrete realisation of these oracles as commitments follows the BCS
transformation of \cref{sec:bcs}. For a polynomial of length $8$ with
blowup $2$, so $N = 16$, the process runs $\log_2 N = 3$ rounds down to
a constant final layer.


\subsection{Query Phase}
\label{subsec:query_phase}

Committing to the codewords is not enough on its own: the verifier must
check that consecutive layers are actually related by the fold. It does
so by sampling a small number of random positions and, at each one,
asking the prover for just the values needed to recompute one fold.

For a single query the verifier picks an index and follows it down
through the layers. At round $r$, with domain size $n_r$, the index
determines a pair of positions half the list apart,
%
\[
\mathrm{lo} = \mathrm{idx} \bmod (n_r/2),
\qquad
\mathrm{hi} = \mathrm{lo} + n_r/2 ,
\]
%
which are exactly the entries $f(x)$ and $f(-x)$ that
\cref{eq:fold_evaluations} combines. The prover returns both values
together with their Merkle authentication paths. The index for the next
round is then $\mathrm{idx} \leftarrow \mathrm{lo}$, so one random choice
in the first layer determines the whole chain of positions down to the
final layer.

In IOP terms, the verifier samples a number of indices in $[0, N/2)$
and, for each one, queries the two oracle positions per layer needed to
recompute one fold. Each query thus costs one pair of oracle queries per
round.


\subsection{Verification and the Colinearity Check}
\label{subsec:verification}

The verifier holds only the proof, namely the Merkle roots, the final
layer, and the opened values with their paths, and must decide whether to
accept. It proceeds in three steps.

First, it \textbf{reconstructs the challenges}. Re-reading the roots and
the final layer from the transcript in the same order the prover wrote
them, it recomputes every $\alpha_r$ and every domain parameter, and
checks that the number of rounds matches the expected
$\log_2(N/\ell)$.

Second, for each queried position it \textbf{checks the Merkle paths},
confirming that the two opened values really sit at their claimed
positions under the committed root. A failed path means the prover is
opening something it did not commit to.

Third, it performs the \textbf{colinearity check}. Using the opened pair
$(\mathrm{lo}, \mathrm{hi})$ at round $r$, it recomputes the folded value
by the same rule the prover should have used,
%
\begin{equation}
\label{eq:colinearity}
v = \frac{\mathrm{lo} + \mathrm{hi}}{2}
  + \alpha_r \cdot \frac{\mathrm{lo} - \mathrm{hi}}{2\,x},
\qquad x = \mathrm{offset}_r \cdot \omega_r^{\,\mathrm{lo}},
\end{equation}
%
and compares $v$ against the value at the corresponding position of the
\emph{next} layer, which is an opened entry for intermediate rounds, or
the final layer for the last round. If the two disagree, the fold was not
performed honestly and the verifier rejects.

The name \emph{colinearity check} comes from the geometry of
\cref{eq:colinearity}. The two points $(x, \mathrm{lo})$ and
$(-x, \mathrm{hi})$, together with the folded value at $\alpha_r$, must
lie on a single line, since the line through $f(x)$ and $f(-x)$ evaluated
at $\alpha_r$ is exactly $f^{(\alpha_r)}(x^2)$. Each query verifies one
such local relation per layer; a codeword far from low degree cannot
satisfy all of them at once, so a handful of random queries suffices to
catch a cheating prover.

In IOP terms the verifier reconstructs the challenges from the
transcript, checks that each queried oracle answer is consistent (which,
after compilation, means checking a commitment opening), and verifies the
colinearity relation \cref{eq:colinearity} at every queried position of
every layer.

\section{Soundness}
\label{sec:fri_soundness}

The earlier sections describe an honest prover. Soundness asks the
opposite: if the committed function is \emph{not} close to any low-degree
polynomial, can a cheating prover still make the verifier accept? We show
the answer is no, except with small probability, following the
round-by-round analysis of Garreta, Mohnblatt, and Wagner~\cite{GMW25},
which gives the simplest known soundness proof for FRI in the form
implemented here.

One point frames everything that follows. FRI certifies \emph{proximity},
not equality (\cref{rem:proximity_not_membership}): the verifier never
learns that the committed vector is exactly a codeword, only that it is
close to one. Soundness must therefore be stated in terms of distance. We
fix a distance parameter $\delta$ and ask that, if the initial function
is $\delta$-far from the code $\RS[\F, D, \rho]$ in relative Hamming
distance (\cref{def:hamming}), then the verifier rejects with high
probability. The difficulty is that the protocol only ever inspects a few
points, so we must argue that being far from the code cannot quietly
disappear during folding, and that it is caught by the queries.


\subsection{The Two Things That Can Go Wrong}
\label{subsec:two_risks}

A cheating prover has exactly two ways to succeed, and the proof bounds
each in turn.

The first is that \emph{folding rescues a far function}. Each round
replaces the current layer $f_i$ by
$f_{i+1} = f_{i,e} + \alpha_i f_{i,o}$, a random linear combination of its
two halves. If some challenge $\alpha_i$ made $f_{i+1}$ close to the code
even though $f_i$ was far, the prover would have folded its way from a
dishonest layer to an honest-looking one, and no amount of querying would
detect it. So we need folding to \emph{preserve distance}: for all but a
few unlucky $\alpha_i$, a far layer stays far.

The second is that \emph{the queries all miss}. Even if every layer stays
far, the verifier only checks a few random positions. If the far parts of
the final layer happened to avoid every queried position, the cheat would
slip through. So we need enough queries that missing all of them is
unlikely.

The next two subsections address these in order: first the algebraic fact
that makes folding preserve distance, then the theorem that combines it
with the query bound.


\subsection{Why Folding Preserves Distance}
\label{subsec:mca}

Distance preservation is not obvious, precisely because $f_{i+1}$ is a
chosen linear combination of $f_{i,e}$ and $f_{i,o}$, and one might fear a
lucky $\alpha_i$ that makes it close to the code by accident. The property
that rules this out is \emph{correlated agreement}, introduced in the
study of proximity gaps~\cite{BCIKS20}: if a random combination of several
functions is close to the code, then those functions are not just
individually close, they agree with codewords on a \emph{common} set of
positions.

The round-by-round proof uses a sharper version, \emph{mutual} correlated
agreement, introduced in the analysis of WHIR~\cite{WHIR24}. It requires
that the set of positions where the combination agrees with the code is
the \emph{same} set on which the two halves agree, not merely a related
one. This equality of sets is what lets the argument repeat cleanly at
every round: the agreement pattern of one layer passes to the next
without drift, so distance is tracked faithfully all the way down. This is
why we use mutual correlated agreement rather than the older notion.

Crucially, mutual correlated agreement is a \emph{theorem}, not an
assumption, as long as the distance stays below the Johnson radius
$1 - \sqrt{\rho}$ (\cref{thm:johnson})~\cite{GMW25}. This is the regime in
which the soundness of our implementation is established, and we return to
what lies beyond it at the end of the section.


\subsection{The Soundness Theorem}
\label{subsec:rbr}

With distance preservation in hand, the two risks of
\cref{subsec:two_risks} become two error terms.

\begin{theorem}[Round-by-round soundness of FRI, informal
  {\cite[Theorem~5.2]{GMW25}}]
\label{thm:rbr}
Let $f$ be $\delta$-far from $\RS[\F, D, \rho]$ with $\delta$ below the
Johnson radius $1 - \sqrt{\rho}$. Then any cheating prover makes the
verifier accept with probability at most
%
\[
\underbrace{\frac{r}{|\mathbb{K}|}\cdot\mathrm{poly}(|D|)}
  _{\text{folding}}
\;+\;
\underbrace{(1 - \delta)^{s}}
  _{\text{queries}} ,
\]
%
where $\mathbb{K}$ is the field the challenges are drawn from,
$r = \log_2(|D|/\ell)$ the number of rounds, and $s$ the number of
queries.
\end{theorem}

\begin{proof}[Proof sketch]
The proof tracks one bit of information through the protocol: whether the
current layer is still far from its code. Following~\cite{GMW25} we call
the state \emph{doomed} while this holds. The argument is that a doomed
state stays doomed through folding, and that a doomed final state is
caught by the queries.

\emph{The start is doomed.} By assumption $f_0$ is $\delta$-far from the
code, so the state is doomed before any interaction. The prover has
committed to a far vector and cannot take it back.

\emph{Folding keeps it doomed} (\cite[Lemma~5.1]{GMW25}). Suppose layer
$f_i$ is $\delta$-far. Its fold is
$f_{i+1} = f_{i,e} + \alpha_i f_{i,o}$. If $f_{i+1}$ were close to the
code for many values of $\alpha_i$, mutual correlated agreement
(\cref{subsec:mca}) would force $f_{i,e}$ and $f_{i,o}$ to agree with
codewords on a common set large enough to make $f_i$ itself close, which
contradicts the assumption. Hence only a few $\alpha_i$ are bad, a
fraction of about $1/|\mathbb{K}|$ of the field. A union bound over the
$r$ rounds gives the folding term. This is the only place the field size
enters, and it is why $\mathbb{K}$ must be the large extension $\Fpk$: a
challenge from the base field would make $1/|\mathbb{K}|$ far too large.

\emph{The queries catch a doomed final layer} (\cite[Claim~5.3]{GMW25}).
If every fold kept the state doomed, the final layer sent in the clear is
itself $\delta$-far from a low-degree polynomial. A single query picks a
random position and checks the colinearity relation \cref{eq:colinearity}
along the chain; since the final layer disagrees with the code on at least
a $\delta$ fraction of positions, the query exposes a broken fold with
probability at least $\delta$. The $s$ queries are independent, so all of
them miss with probability at most $(1-\delta)^s$, the query term.

The prover wins only if some fold got lucky or every query missed, so the
total is the sum of the two terms.
\end{proof}

The two terms are controlled by different means, which is what makes the
protocol practical. The folding term is negligible for free: with
$\mathbb{K} = \Fpk$ of size about $2^{124}$ (\cref{rem:ext_field_fold})
and only a handful of rounds, it sits around $2^{-122}$. The query term is
the one the designer tunes, by choosing how many queries to make.

\begin{example}
\label{ex:query_soundness}
Take a final layer that disagrees with the code on half its positions,
$\delta = 1/2$, and $s = 20$ queries. One query misses with probability
$1/2$, so twenty independent queries all miss with probability
$(1/2)^{20} \approx 10^{-6}$, about $20$ bits of security from the query
phase. A smaller rate raises the detection probability per query and so
buys more bits each. The concrete parameters used in this work, and the
resulting security level, are discussed in \cref{ch:implementation}.
\end{example}


\subsection{Beyond the Provable Regime}
\label{subsec:regimes}

Everything above lives below the Johnson radius $1 - \sqrt{\rho}$, where
mutual correlated agreement is proved. One can push further, up to the
list-decoding \emph{capacity} $1 - \rho$, by \emph{conjecturing} that the
same distance preservation still holds; this would allow fewer queries for
the same security. That regime is attractive for efficiency but unproven,
and recent work has cast doubt on the strongest such conjectures. We
therefore rely only on the provable Johnson-radius guarantee and leave the
capacity regime out of scope.